\chapter{Discuss\~ao}
\label{cap:discussao}

% TODO: capitulo pendente -- redigir apos a execucao empirica prevista no
% cronograma (Anexo I). Estrutura sugerida pelo autor:
%   6.1. Teste das hipoteses H1, H2 e H3 a luz dos resultados
%   6.2. Dialogo com o referencial teorico (assimetria de informacao,
%        capacidade estatal, RAG, XAI)
%   6.3. Limitacoes do estudo (ex.: Sharpe ratio em janela de 12 meses)

\textit{[Cap\'itulo pendente -- a ser redigido ap\'os a execu\c{c}\~ao
emp\'irica prevista no cronograma (Anexo~\ref{ane:cronograma}). Estrutura
sugerida:]}

\begin{itemize}
  \item[] \textit{6.1. Teste das hip\'oteses H1, H2 e H3 \`a luz dos resultados}
  \item[] \textit{6.2. Di\'alogo com o referencial te\'orico (assimetria de informa\c{c}\~ao, capacidade estatal, RAG, XAI)}
  \item[] \textit{6.3. Limita\c{c}\~oes do estudo (ex.: Sharpe ratio em janela de 12 meses)}
\end{itemize}
